Retrieval-Augmented Generation (RAG) has emerged as a practical paradigm for improving downstream tasks by retrieving relevant knowledge from external data sources. It has been successfully deployed in a wide range of real-world applications, including healthcare~\citep{xu2024ram}, law~\citep{wiratunga2024cbr}, finance~\citep{zhang2023enhancing}, and education~\citep{miladi2024leveraging}. With the advent of Large Language Models (LLMs), integrating retrieval into generation further improves faithfulness by mitigating hallucinations and enhancing robustness~\cite{zhao2023survey, huang2023survey}. In most existing RAG systems, retrieval is performed over text corpora.

Graphs provide an explicit representation of relational structure and have long been used across domains such as knowledge representation, social networks, and biomedical discovery~\cite{wu2020comprehensive, ma2021deep, wu2023survey}. Graph Retrieval-Augmented Generation (GraphRAG) has recently gained attention for retrieving and aggregating information from graph-structured data, including knowledge graphs (KGs) and molecular graphs~\citep{han2024retrieval, peng2024graph}. Beyond leveraging existing graphs, an emerging line of work extends GraphRAG to text-based tasks by constructing graphs from unstructured documents, with reported benefits for global summarization~\cite{edge2024local}, planning~\cite{lin2024graph}, and reasoning~\cite{han2025reasoning}. However, most GraphRAG studies for text are conducted under task- and system-specific settings, using bespoke datasets, graph construction heuristics, and evaluation protocols. This heterogeneity makes it difficult to draw principled conclusions about \textit{when and why} explicit graph structures help (or hurt) retrieval-augmented generation, and obscures practical trade-offs such as construction cost, retrieval latency, and storage footprint.

To address this gap, we conduct a controlled and systematic benchmark of RAG and GraphRAG on widely used text-based tasks, focusing on \textit{Question Answering} (QA) and \textit{Query-based Summarization}. We consider four representative categories of GraphRAG systems: 
{\bf (1)} \emph{KG-based GraphRAG}~\cite{Liu_LlamaIndex_2022}, which extracts a KG from text and performs retrieval over the KG; 
{\bf (2)} \emph{Community-based GraphRAG}~\citep{edge2024local}, which performs retrieval over community structures and hierarchical abstractions; 
{\bf (3)} \emph{Text-centric graph-guided RAG}~\citep{jimenez2024hipporag}, which retrieves original text chunks with the assistance of a constructed knowledge graph; and 
{\bf (4)} \emph{Hierarchical summary-based GraphRAG}~\citep{sarthi2024raptor}, which builds hierarchical summaries to enable multi-granular retrieval without relying on explicit KGs.
For QA, we evaluate both single-hop and multi-hop settings; for summarization, we consider both single-document and multi-document scenarios. Crucially, we introduce a unified evaluation protocol that standardizes data preprocessing, retrieval, and generation settings, enabling fair and reproducible comparisons across paradigms.

Our analysis results lead to several key findings. First, \textbf{RAG and GraphRAG exhibit complementary behaviors rather than a consistent winner}. In QA, RAG performs better on single-hop and detail-oriented factual queries, whereas GraphRAG is more effective on multi-hop, reasoning-intensive questions. Second, \textbf{GraphRAG design choices matter}: for example, community-based global search can sacrifice query-specific details, hurting detail-oriented QA, while providing more corpus-level aggregation that benefits broad or diverse summarization outputs. Third, \textbf{evaluation protocol can change conclusions}: we show that LLM-as-a-Judge evaluation for summarization can be highly sensitive to the presentation order of candidate summaries, introducing strong position effects that may confound comparisons. Finally, \textbf{GraphRAG is not free}: it often incurs higher construction cost, retrieval latency, and storage footprint, and its performance can be sensitive to the quality (and cost) of graph construction.
These findings suggest that effective retrieval-augmented generation should not treat RAG and GraphRAG as mutually exclusive choices. Motivated by this, we study two practical hybrid strategies: \textbf{Selection}, which routes queries to RAG or GraphRAG based on query type for efficiency, and \textbf{Integration}, which combines evidence from both paradigms to maximize performance. Across benchmarks, these strategies yield consistent improvements.
Our main contributions are as follows:
\begin{itemize}[leftmargin=*, itemsep=0pt, parsep=0pt]
    \item \textbf{Systematic Benchmark:} We present a controlled benchmark comparing RAG and multiple GraphRAG variants across QA and query-based summarization under a unified evaluation protocol (consistent preprocessing, retrieval, and generation settings) for fair and reproducible comparison.
    \item \textbf{Strong, Task-Level Findings:} We identify clear complementarities: RAG is stronger for factual/detail-oriented QA, while GraphRAG benefits reasoning-intensive QA and produces more corpus-level, diverse summaries, with outcomes strongly affected by GraphRAG design choices (e.g., local vs.\ global search).
    \item \textbf{Hybrid Strategies:} We study \textbf{Selection} and \textbf{Integration} strategies that combine RAG and GraphRAG, achieving consistent improvements and illustrating effectiveness--efficiency trade-offs.
    \item \textbf{Evaluation and Efficiency Analyses:} We analyze failure modes, construction/retrieval/storage costs, sensitivity to graph construction quality, and demonstrate strong position effects in LLM-as-a-Judge summarization evaluation, highlighting practical considerations for reliable (Graph)RAG assessment.
\end{itemize}


% Our results indicate that RAG and GraphRAG exhibit \emph{complementary behaviors} rather than a consistent dominance of one over the other. For QA, RAG tends to perform better on single-hop and detail-oriented queries, whereas GraphRAG is more effective on reasoning-intensive, multi-hop questions. For query-based summarization, RAG typically preserves fine-grained, query-specific details, while GraphRAG often produces more diverse and corpus-level summaries. We further show that commonly used LLM-as-a-Judge evaluation for summarization can introduce strong \emph{position effects}, which may substantially affect conclusions. Finally, we analyze practical considerations, including indexing time, retrieval latency, generation cost, and token/storage usage, and show that GraphRAG performance can be sensitive to graph construction quality. Motivated by these findings, we explore two hybrid strategies, \textbf{Selection} and \textbf{Integration}, to combine the strengths of RAG and GraphRAG and improve overall performance.
% Our main contributions are as follows:
% % \vspace{-0.05in}
% \begin{itemize}[leftmargin=*, itemsep=0pt, parsep=0pt]
%     \item \textbf{Systematic Evaluation:} We present a controlled benchmark comparing RAG and representative GraphRAG variants across question answering and query-based summarization under a unified protocol, enabling fair and reproducible evaluation.
%     \item \textbf{Task-Specific Insights:} We provide an in-depth analysis of when RAG vs.\ GraphRAG is preferable, highlighting complementary strengths across query types and task objectives.
%     \item \textbf{Hybrid Retrieval Strategies:} We study two practical hybrid approaches, \textbf{Selection} and \textbf{Integration}, that combine RAG and GraphRAG to leverage their complementary strengths and consistently improve performance.
%     \item \textbf{Evaluation and Efficiency Analyses:} We identify key limitations and practical trade-offs, including position effects in LLM-as-a-Judge summarization evaluation, as well as construction/retrieval/storage costs and sensitivity to graph construction quality, and discuss implications for building more reliable and efficient (Graph)RAG systems.
% \end{itemize}
